\documentclass[12pt]{article}
\usepackage{amsmath,amssymb,amsfonts}
\usepackage[margin=1in]{geometry}

\begin{document}

\begin{center}
{\Large\bfseries Chapter 10, Problem 7}
\end{center}

\noindent\textbf{Problem.} Let $\{123 \cdots l\}$ be a cycle of length $l$. Prove that if $\lambda < l$ then $\{123 \cdots l\}^\lambda \neq e$, while $\{123 \cdots l\}^l = e$.

\bigskip
\noindent\textbf{Solution.}

Let $\sigma = (1\;2\;3\;\cdots\;l)$ be the $l$-cycle. This permutation acts as $\sigma(k) = k+1$ for $k = 1, \ldots, l-1$ and $\sigma(l) = 1$. In general, $\sigma(k) = k + 1 \pmod{l}$ (with representatives in $\{1, 2, \ldots, l\}$).

\medskip
\noindent\textbf{Claim 1:} $\sigma^\lambda(k) = k + \lambda \pmod{l}$ for all $k$.

\textit{Proof by induction on $\lambda$:}

Base case: $\sigma^1(k) = k + 1 \pmod{l}$. \checkmark

Inductive step: Assume $\sigma^\lambda(k) = k + \lambda \pmod{l}$. Then
\[
\sigma^{\lambda+1}(k) = \sigma(\sigma^\lambda(k)) = \sigma(k + \lambda \pmod{l}) = (k + \lambda) + 1 \pmod{l} = k + (\lambda + 1) \pmod{l}. \; \checkmark
\]

\medskip
\noindent\textbf{Claim 2:} $\sigma^l = e$.

Using Claim 1: $\sigma^l(k) = k + l \equiv k \pmod{l}$ for all $k \in \{1, \ldots, l\}$. Therefore $\sigma^l$ is the identity. $\square$

\medskip
\noindent\textbf{Claim 3:} If $1 \leq \lambda < l$, then $\sigma^\lambda \neq e$.

Using Claim 1: $\sigma^\lambda(1) = 1 + \lambda \pmod{l}$. Since $1 \leq \lambda < l$, we have $1 + \lambda \not\equiv 1 \pmod{l}$ (because $\lambda \not\equiv 0 \pmod{l}$). Therefore $\sigma^\lambda(1) \neq 1$, which means $\sigma^\lambda \neq e$. $\square$

\medskip
Together, these results show that the $l$-cycle $(1\;2\;3\;\cdots\;l)$ has order exactly $l$ in the symmetric group.

\end{document}
